\chapter{微分表达式、算子域与自伴性}
\label{chap:math-unbounded}

在区间 \((0,1)\) 上，\(-f''\) 配固定端点条件有正本征值 \(n^2\pi^2\)；配自由端点条件则多出零本征值。两个问题使用同一个微分表达式，却给出不同的可测能量。算子必须同时指定它作用的空间、作用规则和允许输入的函数。

\section{定义域为什么不能省略}

\begin{definition}[稠密定义的算子与伴随]
在复 Hilbert 空间 \(\mathcal H\) 中，线性算子 \(A\) 是定义域 \(D(A)\subset\mathcal H\) 和线性规则 \(A:D(A)\to\mathcal H\) 的组合。若 \(D(A)\) 的闭包等于 \(\mathcal H\)，称它稠密定义。此时
\[
D(A^*)=\{y\in\mathcal H:\text{存在 }z\in\mathcal H\text{ 使 }\langle y,Ax\rangle=\langle z,x\rangle\text{ 对一切 }x\in D(A)\},
\]
并定义 \(A^*y=z\)。稠密性保证这样的 \(z\) 唯一。
\end{definition}

此定义要求 \(x\mapsto\langle y,Ax\rangle\) 在原空间范数下有界；它不仅要求形式积分分部。若 \(D(A)\) 不稠密，则 \(z\) 可以加上任意 \(D(A)^\perp\) 而不改变等式，伴随不再是单值算子。

微分算子的自然域要允许弱导数。区间上的 \(H^1(0,1)\) 指 \(f\in L^2\) 且存在 \(g\in L^2\) 使 \(\int f\varphi'=-\int g\varphi\) 对一切紧支光滑 \(\varphi\) 成立，记 \(g=f'\)。\(H^2\) 再要求 \(f'\in H^1\)。一维 \(H^1\) 函数可取绝对连续代表，因而端点值有意义；\(H_0^1\) 则是两个端点值都为零的 \(H^1\) 函数。

这里的“端点值”并非任意给一个等价类挑代表。若 \(f\in H^1(0,1)\)，以 \(g=f'\) 作普通积分，令 \(F(x)=\int_0^x g(s)\,\dd s\)。按分布导数的定义，\((f-F)'=0\)。导数为零的分布在连通区间上必为常数：任取积分为零的紧支测试函数 \(\varphi\)，其原函数 \(\Phi(x)=\int_{-\infty}^x\varphi(s)\,\dd s\) 仍紧支，故 \(\langle f-F,\varphi\rangle=-\langle(f-F)',\Phi\rangle=0\)；任意测试函数减去其积分乘一支固定的单位积分测试函数后都属于这一类。因此 \(f(x)=c+F(x)\) 几乎处处，并且右端连续，正是唯一的连续代表。由 Cauchy--Schwarz，
\[
 |f(x)-f(y)|\le\|f'\|_2|x-y|^{1/2},\qquad
 |f(0)|\le \|f\|_2+\frac1{\sqrt3}\|f'\|_2.
\]
第二个估计把 \(f(0)=\int_0^1f(x)\,\dd x-\int_0^1(1-s)f'(s)\,\dd s\) 代入即可；右端点同理。这证明端点取值是 \(H^1\) 图范数下的连续运算，却不是 \(L^2\) 范数下的连续运算。边界条件因此能够定义 \(H^1\) 或 \(H^2\) 中的闭子空间。

\(H^1\) 本身对 \(\|f\|_{H^1}^2=\|f\|_2^2+\|f'\|_2^2\) 完备。若 \((f_n)\) 对此范数为 Cauchy 列，\(L^2\) 的完备性给出 \(f_n\to f\)、\(f_n'\to g\)。把 \(\int f_n\varphi'=-\int f_n'\varphi\) 对固定测试函数取极限，就得到 \(f'=g\) 的弱导数关系。连续的端点取值保证零端点条件在极限下仍成立，所以 \(H_0^1\) 也完备。这个小证明在后文会同时保证动量算子闭合以及能量形式闭合。

取 \(\mathcal H=L^2(0,1)\)，先在 \(C_c^\infty(0,1)\) 上定义 \(P_0f=-\ii f'\)。对 \(f\in H^1\) 与测试函数 \(g\) 分部积分，\(\langle f,P_0g\rangle=\langle-\ii f',g\rangle\)。反过来，若 \(f\in D(P_0^*)\)，伴随定义对所有测试函数成立，正是 \(f\) 的弱导数属于 \(L^2\) 的条件。因此
\[
D(P_0^*)=H^1(0,1),\qquad P_0^*f=-\ii f'.
\]
同一个作用规则已经出现两个不同的域；端点条件将在自伴性中决定下来。

\begin{definition}[闭算子]
\(A\) 的图为 \(\mathcal G(A)=\{(x,Ax):x\in D(A)\}\subset\mathcal H\oplus\mathcal H\)。若 \(x_n\to x\) 且 \(Ax_n\to y\) 总推出 \(x\in D(A)\)、\(Ax=y\)，称 \(A\) 闭。若图的闭包仍是一个算子的图，称 \(A\) 可闭，该算子记为 \(\overline A\)。
\end{definition}

闭性不是定义域本身闭。导数算子 \(P_0\) 在 \(C_c^\infty\) 上不闭：光滑紧支函数可在 \(H^1\) 范数下逼近有零端点值而非紧支的函数。它的闭包为 \((-\ii\partial_x,H_0^1(0,1))\)。图范数 \(\|f\|_A^2=\|f\|^2+\|Af\|^2\) 恰好同时控制输入与输出；闭性等价于 \(D(A)\) 对图范数完备。

\section{分部积分能证明什么}

区间动量算子的闭包 \(P=(-\ii\partial_x,H_0^1(0,1))\) 对称。对 \(f,g\in H^1(0,1)\)，逐次使用一维分部积分得
\[
\langle f,-\ii g'\rangle-\langle-\ii f',g\rangle
=-\ii[\overline f g]_0^1.
\]
若两函数端点均为零，边界项消失，所以 \(P\subset P^*\)。然而已算出 \(D(P^*)=H^1(0,1)\)，严格大于 \(H_0^1(0,1)\)：常数函数就在伴随域，却不在 \(D(P)\)。这给出一个对称而不自伴的完整反例。

\begin{definition}[对称与自伴]
稠密定义算子 \(A\) 对称，意为 \(\langle f,Ag\rangle=\langle Af,g\rangle\) 对所有 \(f,g\in D(A)\) 成立，即 \(A\subset A^*\)。若还满足 \(D(A)=D(A^*)\)，从而 \(A=A^*\)，称 \(A\) 自伴。
\end{definition}

端点条件可以补足动量算子的定义域。固定 \(\theta\in[0,2\pi)\)，令
\[
D(P_\theta)=\{f\in H^1(0,1):f(1)=\ee^{\ii\theta}f(0)\},
\qquad P_\theta f=-\ii f'.
\]
该域上边界项为零。要检查伴随域，不能止于这句话：若 \(g\in D(P_\theta^*)\)，先以紧支函数测试得 \(g\in H^1\)；再让 \(f(0)\) 任取复数，且 \(f(1)=\ee^{\ii\theta}f(0)\)，边界项消失要求 \(\overline{g(1)}\ee^{\ii\theta}=\overline{g(0)}\)，即 \(g(1)=\ee^{\ii\theta}g(0)\)。故 \(D(P_\theta^*)=D(P_\theta)\)，它确实自伴。解 \(-\ii f'=\lambda f\) 得 \(f(x)=C\ee^{\ii\lambda x}\)，边界条件给 \(\lambda=2\pi n+\theta\)，\(n\in\mathbb Z\)。

\section{二阶算子的边界型}

对 \(S_0=-\partial_x^2\) 取初始域 \(C_c^\infty(0,1)\)。若 \(g\in D(S_0^*)\)，则某个 \(h\in L^2\) 满足 \(\int\overline g(-\varphi'')=\int\overline h\varphi\) 对所有紧支光滑 \(\varphi\) 成立。这正说 \(-g''=h\) 为弱导数。把 \(h\) 积分两次，并减去所得原函数，余下的分布二阶导数为零，因而是仿射函数；所以 \(g\in H^2(0,1)\)。反向分部积分立即成立，故 \(D(S_0^*)=H^2(0,1)\)，称最大域。\(S_0\) 的图闭包的域为 \(H_0^2(0,1)\)，即 \(H^2\) 中函数值和一阶导数在两端均为零的函数；这称最小域。注意最小域比 Dirichlet 域还小，因为后者不限制端点导数。在最大域中，分部积分两次给出
\[
\langle f,S_0^*g\rangle-\langle S_0^*f,g\rangle
=\bigl[\overline{f'}g-\overline f g'\bigr]_0^1.
\]
这个边界型同时含函数值和导数值，所以二阶算子比动量多需要一组端点条件。若取 \(D(S_D)=\{f\in H^2:f(0)=f(1)=0\}\)，边界型在域内消失；若 \(g\in D(S_D^*)\)，令 \(f'(0),f'(1)\) 独立变化，又迫使 \(g(0)=g(1)=0\)。于是 \(S_D\) 自伴。固定导数为零的 Neumann 域与同样推理得到自伴；它容许常数零模，而 Dirichlet 域不容许。

独立改变端点条件不能只看边界型的一个方向。比如只要求 \(f(0)=f(1)=0\) 而忘记域必须落在 \(H^2\)，那么 \(-f''\) 可能根本不属于 \(L^2\)，算子尚未定义。完整检查总是先确定最大域，再要求边界型在候选域内为零，最后算候选域的伴随是否仍是它自己。

Robin 条件也可把边界响应连续地反映到能谱。两端取同一非负参数 \(h\)：\(f'(0)=hf(0)\)、\(f'(1)=-hf(1)\)。设正本征值为 \(k^2\)，通解 \(f=A\cos kx+B\sin kx\)。左条件给 \(B=hA/k\)；把它代入右条件并消去非零 \(A\)，得到
\[
(h^2-k^2)\sin k+2hk\cos k=0.
\]
当 \(h=0\) 时还要单独检查 \(k=0\)，它给 Neumann 零模；当 \(h>0\) 时零模消失。这个小计算展示边界参数进入谱的具体位置，也提醒我们除以 \(k\) 前必须先处理 \(k=0\)。

\begin{exercise}
在 \(D(S_R)=\{f\in H^2(0,1):f'(0)=\alpha f(0),\ f'(1)=-\beta f(1)\}\) 中取实数 \(\alpha,\beta\)。用边界型检验对称性，再通过任意选择 \(f(0),f(1)\) 求 \(D(S_R^*)\)，证明它自伴。
\end{exercise}

\section{亏量空间记录缺少的边界数据}

区间动量在零端点域上对称却不自伴；给端点增加相位关系以后成为自伴。能否一般地判断一个对称算子是否容许这样的补充？非实数 \(\pm\ii\) 不可能是对称算子的本征值，因为对称算子的本征值实。因此伴随在这两个值上的核只记录原定义域尚未控制的方向。

\begin{definition}[亏量空间]
对稠密定义、闭且对称的算子 \(S\)，定义 \(\mathcal N_+=\ker(S^*-\ii I)\)、\(\mathcal N_-=\ker(S^*+\ii I)\)，其维数 \(n_+,n_-\) 称亏格。
\end{definition}

在区间动量 \(P=(-\ii\partial_x,H_0^1(0,1))\) 上，\(P^*f=\ii f\) 给 \(f'=-f\)，所以 \(\mathcal N_+=\operatorname{span}\{\ee^{-x}\}\)；\(P^*f=-\ii f\) 给 \(f'=f\)，所以 \(\mathcal N_-=\operatorname{span}\{\ee^x\}\)。有限区间上二者均平方可积，亏格为 \((1,1)\)。端点相位 \(\theta\) 正是一维亏空间间酉映射的一个复相位。

\begin{theorem}[von Neumann 扩张定理]
设 \(S\) 为稠密定义、闭且对称的算子。它有自伴扩张当且仅当 \(n_+=n_-\)。每个自伴扩张由酉映射 \(U:\mathcal N_+\to\mathcal N_-\) 给出，其域为
\[
D(S_U)=D(S)\mathbin{\dot+}\{u+Uu:u\in\mathcal N_+\},
\qquad S_U(f+u+Uu)=Sf+\ii u-\ii Uu.
\]
\end{theorem}

这里 \(\dot+\) 表示每个向量具有唯一的分解。为了看清为什么两个亏格必须相等，不妨把证明中的关键线性代数真正展开。伴随算子 \(S^*\) 是闭的，所以在 \(D(S^*)\) 上，
\[
 (f,g)_G=\langle f,g\rangle+\langle S^*f,S^*g\rangle
\]
给出一个完备的内积。闭性也使 \(D(S)\) 成为这个图内积下的闭子空间。令 \(M\) 是它的正交补。若 \(h\in M\)，则对一切 \(f\in D(S)\) 有
\[
 \langle h,f\rangle+\langle S^*h,Sf\rangle=0.
\]
后一等式按伴随的定义说明 \(S^*h\in D(S^*)\)，而且 \(S^*(S^*h)=-h\)。它还说明 \(S^*h\in M\)：把 \(h\) 换成 \(S^*h\) 再与任意 \(f\in D(S)\) 取图内积，所得两项恰好相消。因此 \(S^*\) 在 \(M\) 上的平方是 \(-I\)。任取 \(h\in M\)，直接计算
\[
 h_+=\frac{h-\ii S^*h}{2},\qquad
 h_-=\frac{h+\ii S^*h}{2},\qquad
 S^*h_+=\ii h_+,\quad S^*h_-=-\ii h_-.
\]
我们得到伴随域的分解
\begin{equation}
 D(S^*)=D(S)\mathbin{\dot+}\mathcal N_+\mathbin{\dot+}\mathcal N_-.
 \label{eq:math-von-neumann-decomposition}
\end{equation}
三项在图内积下两两正交；例如 \(u\in\mathcal N_+\)、\(v\in\mathcal N_-\) 时，\((u,v)_G=\langle u,v\rangle+\langle\ii u,-\ii v\rangle=0\)。这一步解释了为何原算子域之外恰好只有两个“未配对”的方向，而不是任意多种边界数据。

现在把对称性写成伴随域上的边界型
\[
 B(f,g)=\langle f,S^*g\rangle-\langle S^*f,g\rangle.
\]
含有 \(D(S)\) 分量的项全为零。对 \(u_+,v_+\in\mathcal N_+\) 与 \(u_-,v_-\in\mathcal N_-\)，直接代入 \(S^*u_+=\ii u_+\)、\(S^*u_-=-\ii u_-\) 得
\[
 B(u_++u_-,v_++v_-)
 =2\ii\bigl(\langle u_+,v_+\rangle-\langle u_-,v_-\rangle\bigr).
\]
两种亏方向的交叉项自动为零。若把扩张域选成 \(D(S)\dot+\{u+Uu:u\in\mathcal N_+\}\)，上式为零的充要条件是 \(U\) 保内积。只有当 \(U\) 还映满 \(\mathcal N_-\) 时，扩张的伴随域才不会留下额外方向。确实，设 \(g=f+v_++v_-\in D(S_U^*)\)，逐一拿 \(u+Uu\) 测试边界型，得到 \(\langle v_+,u\rangle=\langle v_-,Uu\rangle\) 对一切 \(u\)。若 \(U\) 酉，则 \(v_-=Uv_+\)，故 \(g\in D(S_U)\)。

反向也能直接检验，而不必把“极大”当作未解释的口号。设 \(A\) 是 \(S\) 的任意自伴扩张；其域在分解~\cref{eq:math-von-neumann-decomposition} 中除 \(D(S)\) 外给出某个子空间 \(L\subset\mathcal N_+\oplus\mathcal N_-\)。对 \((u_+,u_-)\in L\) 取 \(B(u_++u_-,u_++u_-)=0\)，得到 \(\|u_+\|=\|u_-\|\)。所以投影 \(L\to\mathcal N_+\) 是单射，并把 \(L\) 写成一个保内积映射 \(U:M\to\mathcal N_-\) 的图，其中 \(M\subset\mathcal N_+\)。由于自伴算子闭，\(L\) 对图范数闭；在 \(L\) 上两个投影的范数相等，因而 \(M\) 也是闭子空间。若 \(M\ne\mathcal N_+\)，便可取非零 \(v_+\perp M\)。边界型公式给 \(B(v_+,u+Uu)=0\) 对所有 \(u\in M\)，所以 \(v_+\in D(A^*)\)；但 \(v_+\notin D(A)\)，与自伴性矛盾。故 \(M=\mathcal N_+\)。若 \(U(M)\ne\mathcal N_-\)，同样取非零 \(v_-\perp U(M)\)，得到 \(v_-\in D(A^*)\setminus D(A)\)，仍矛盾。于是 \(U\) 必须映满，正是两亏空间间的酉映射；它存在当且仅当两空间维数相等。

图内积的证明并没有假定微分表达式或端点值，所以定理可用于一般闭对称算子；具体端点条件仍须把 \(u+Uu\) 翻译回函数的边界数据。它并不说任意写一个端点条件就能自伴：条件必须让伴随域恰好等于所选域。

在半直线 \((0,\infty)\) 上重新取 \(P=-\ii\partial_x\) 与域 \(H_0^1(0,\infty)\)。伴随域仍为 \(H^1(0,\infty)\)，亏量方程仍给 \(\ee^{-x}\) 与 \(\ee^x\)，但后者不平方可积，于是亏格为 \((1,0)\)。定理表明半直线动量没有自伴扩张。端点只能让平移向一个方向离开半直线，无法形成对所有正负时间都可逆的酉平移群。

同在半直线，二阶表达式却给不同答案。令 \(S=-\partial_x^2\) 从 \(C_c^\infty(0,\infty)\) 闭合而来。伴随最大域中的亏量方程 \(-f''=\pm\ii f\) 各有两支指数解，但每个符号都恰有一支在无穷远平方可积，所以 \((n_+,n_-)=(1,1)\)。一个边界参数足以形成自伴扩张：\(D(S_\gamma)=\{f\in H^2(0,\infty):f'(0)=\gamma f(0)\}\)，其中 \(\gamma\in\mathbb R\)。无穷远边界型对 \(H^2\) 函数消失；在零点，候选域内边界型消失，而让 \(f(0)\) 任取又使伴随域满足同一关系。若 \(\gamma<0\)，\(f(x)=\sqrt{-2\gamma}\,\ee^{\gamma x}\) 归一化且满足 \(S_\gamma f=-\gamma^2f\)。同一自由微分式会因边界吸引产生一个负能束缚态。

\section{同一能量表达式的三种自伴实现}

对 \(-\partial_x^2\) 的 Dirichlet 域，归一化本征函数为 \(\sqrt2\sin(n\pi x)\)，本征值 \(n^2\pi^2\)，\(n\ge1\)。Neumann 域的本征函数为 \(1\) 与 \(\sqrt2\cos(n\pi x)\)，本征值 \(0\) 与 \(n^2\pi^2\)。若两端接起来，采用 \(f(1)=f(0)\)、\(f'(1)=f'(0)\)，则本征函数为 \(\ee^{2\pi\ii nx}\)，本征值 \((2\pi n)^2\)，\(n\in\mathbb Z\)。这三种谱差别完全来自定义域。

另一个连续可调的选择是实参数 \(\alpha,\beta\) 的 Robin 条件 \(f'(0)=\alpha f(0)\)、\(f'(1)=-\beta f(1)\)。它仍令边界型消失，且伴随计算会重得同样两个条件。物理上，参数改变了边界处的响应；数学上，它改变了自伴实现，也就改变能级。闭区间二阶算子还有更一般的耦合端点条件，亏量空间的二维酉映射给出其系统分类。

这些显式本征函数还允许直接核验自伴性与演化的关系。这里区间坐标 \(x\in(0,1)\) 已无量纲化，\(S_D=-\partial_x^2\) 的本征值也是无量纲数；取无量纲演化参数 \(s\)。Fourier 正弦系给 \(\psi=\sum_{n\ge1}c_n\sqrt2\sin(n\pi x)\) 且 \(\sum|c_n|^2=\|\psi\|^2\)。定义
\[
U(s)\psi=\sum_{n\ge1}\ee^{-\ii n^2\pi^2s}c_n\sqrt2\sin(n\pi x).
\]
每个系数只乘模为一的相位，因此 Parseval 恒等式给 \(\|U(s)\psi\|=\|\psi\|\)，而 \(U(s+r)=U(s)U(r)\) 逐系数成立。对满足 \(\sum n^4|c_n|^2<\infty\) 的态，可逐项在 \(L^2\) 中求导并得到 \(\ii\partial_s\psi=S_D\psi\)。若要恢复物理长度 \(L\)、粒子质量 \(m\) 和时间 \(t\)，Hamiltonian 是 \(H=\hbar^2S_D/(2mL^2)\)，所以 \(s=\hbar t/(2mL^2)\)。Stone 定理把这个逐模构造推广到任何自伴算子，包括没有离散本征函数基的情形。

\begin{exercise}
对半直线上的 \(-\partial_x^2\) 与初始域 \(C_c^\infty(0,\infty)\)，解 \(-f''=\pm\ii f\) 并逐个判断平方可积性，得到其亏格。再说明为何零端点值和实 Robin 条件都可能给自伴扩张，而半直线动量不能。
\end{exercise}


\section{自伴性与概率守恒}

若 \(H=H^*\)，谱定理定义 \(U(t)=\exp(-\ii tH/\hbar)\)，且 \(U(t)^*U(t)=I\)。因此任意初态 \(\psi\in\mathcal H\) 都满足 \(\|U(t)\psi\|=\|\psi\|\)，即总概率守恒。对于 \(\psi\in D(H)\)，它还满足强意义的 Schrödinger 方程 \(\ii\hbar\dot\psi=H\psi\)。Stone 定理的逆向说明：每个强连续酉时间演化群都有唯一自伴生成元。这里引用的是一般定理；上面的区间算例则亲手检验了自伴域怎样让它适用。

只说 \(H\) 对称还不够。对称性只允许对域内态形式地计算
\[
\frac{\dd}{\dd t}\|\psi(t)\|^2
=\frac{\ii}{\hbar}\bigl(\langle H\psi,\psi\rangle-\langle\psi,H\psi\rangle\bigr)=0,
\]
却没有保证任意初态都存在双向时间演化。半直线动量的亏格不相等，正好给出这种形式计算失效的例子。下一章要问：既然自伴算子可以生成演化，怎样在没有普通本征向量的情况下计算它？
